% Generated draft: requires mathematical and source review.
\documentclass[11pt]{article}
\usepackage[margin=25mm]{geometry}
\usepackage{fontspec}
\setmainfont{Latin Modern Roman}
\usepackage{amsmath,amssymb,mathtools,unicode-math}
\setmathfont{Latin Modern Math}
\usepackage{enumitem,needspace,xurl,hyperref}
\setlength{\parindent}{0pt}
\setlength{\parskip}{5pt}
\setlength{\emergencystretch}{3em}
\newcommand{\N}{\mathbb{N}}
\newcommand{\Z}{\mathbb{Z}}
\newcommand{\Q}{\mathbb{Q}}
\newcommand{\R}{\mathbb{R}}
\newcommand{\C}{\mathbb{C}}
\newcommand{\F}{\mathbb{F}}
\newcommand{\A}{\mathbb{A}}
\newcommand{\PP}{\mathbb{P}}
\newcommand{\E}{\mathbb{E}}
\newcommand{\eps}{\varepsilon}
\newcommand{\dd}{\,\mathrm{d}}
\begin{document}
\section*{UnsolvedMath ID: 700817}
\textbf{Source problem code:} catalogue entry
\subsection{Definitions and mathematical statement}
Let \(p\) be a prime, let \(G\) be a finite group, and let \(V\leq\operatorname{Aut}(G)\) be an elementary abelian group of order \(p^n\); thus \(V\cong(\Z/p\Z)^n\). For \(v\in V\), define
\[
[G,v]=\langle g^{-1}v(g):g\in G\rangle.
\]
For positive rank \(n\), put
\[
H=\bigcap_{v\in V\setminus\{1\}}[G,v].
\]
The source requires \(V\) to be a \emph{regular} group of automorphisms but does not define this term. A provisional interpretation is
\[
C_G(V):=\{g\in G:v(g)=g\text{ for every }v\in V\}=\{1\}.
\]
This interpretation requires confirmation from the cited source. It must be distinguished from the stronger condition \(C_G(v)=\{1\}\) for every nonidentity \(v\in V\). No assumption that \(p\nmid |G|\) is explicitly supplied.

Subject to confirmation of the meaning of regularity, the first question is: must \(H\) be nilpotent for every such \(G\) and \(V\)? Here, writing \([K,L]=\langle k^{-1}\ell^{-1}k\ell:k\in K,\ell\in L\rangle\), define \(\gamma_1(H)=H\) and \(\gamma_{i+1}(H)=[\gamma_i(H),H]\). Nilpotence means that \(\gamma_{c+1}(H)=\{1\}\) for some integer \(c\geq0\); its least such value is the nilpotency class \(\operatorname{cl}(H)\).

If the first question has an affirmative answer, the source asks whether the nilpotency class of \(H\) is bounded by a function depending only on \(p\), \(n\), and the derived length of \(G\). Define \(G^{(0)}=G\) and \(G^{(i+1)}=[G^{(i)},G^{(i)}]\). When \(G\) is soluble, its derived length \(d(G)\) is the least \(d\geq0\) with \(G^{(d)}=\{1\}\). On this domain the requested bound is a function
\[
f:\{p:p\text{ is prime}\}\times\Z_{\geq1}\times\Z_{\geq0}\longrightarrow\Z_{\geq0}
\]
such that \(\operatorname{cl}(H)\leq f(p,n,d(G))\) for every soluble \(G\) satisfying the hypotheses. The record neither explicitly assumes solubility of \(G\) nor states that it follows from regularity; consequently the intended domain of the second question remains unresolved. The source also does not specify whether rank zero is allowed or how to interpret the empty intersection in that case.
\subsection{Short English statement}
For a finite group admitting a regular elementary abelian automorphism group, intersect the subgroups generated by the commutators with each nonidentity automorphism. Must this intersection be nilpotent, and, if so, can its class be bounded using only the prime, the automorphism group rank, and the derived length of the original group? The meaning of regularity and the solubility qualification need source confirmation.
\subsection{Sources}
\begin{itemize}
\item Local UnsolvedMath registry, supplied complete source record: permanent ID \(700817\), exact source code \texttt{catalogue entry}, title \emph{Kourovka Notebook Problem 11.125}. The registry labels the problem open; this is a supplied status, not an independently verified mathematical or literature assessment.
\item Supplied source attribution: \emph{The Kourovka Notebook (Unsolved Problems in Group Theory)}, original problem 11.125, 11th issue (1990); referenced in the supplied record through the 21st issue edition (2026), \texttt{arXiv:1401.0300v48}. The collection description names E. I. Khukhro and V. D. Mazurov as editors. The cited text has not been independently consulted for this draft.
\item Proposed in 1990. The proposer field is supplied as \texttt{P. V. Shumyatski\textbackslash{}u\{\textbackslash{}i}; its accent suffix is malformed, so the intended complete spelling requires verification.
\end{itemize}
\end{document}
